% P01 — INSERTION into ProofFinal/chapters/60-sns.tex (Section sec:preliminary_sim_results).
% Insert as a new paragraph BETWEEN:
%   end of paragraph: "...reproducing the published behavior of the original model in its original engine."
%   start of next paragraph: "The replication series has a defined extension, and its current boundary is stated plainly...."
% (60-sns.tex at repository commit 39e3f707/d19d4ae2, char offset 50672; see REVIEW.md for the exact boundary text.)
% REV 2 (2026-10-26) after supervisor F1/F2: response attribution split between the
% 2023 biarticular model (RG-stimulus phenomenology) and the antagonist-only
% predecessor (knee-PF no-effect); ports clause scoped to the paper-exact window.

The replication target was then extended from the program's own archived models to the published original project itself. I exported the archived AnimatLab project of the two-layer-CPG rat-hindlimb model of \citet{deng_neuromechanical_2019} (ModelDB entry 264844) to standalone-simulation form with AnimatLab's own serializer, instrumented it with neural voltage charts, and re-ran it unmodified in its original physics engine for \qty{10.1}{\second}. The project free-runs with no external neural drive: each leg's rhythm-generator extensor and flexor populations alternate at a \qty{0.23}{\second} period, and the two legs fire in phase. I then applied the stimulation protocol of \citet{deng_biarticular_2023}, a \qty{2}{\nano\ampere} tonic current on the left rhythm-generator flexor population from \SIrange{2}{2.1}{\second}. Under that stimulus the original shows the same phenomenology that paper reports for its biarticular model under the rhythm-generator stimulus: an additional short flexor burst begins at stimulus onset (\qty{2.000}{\second}, against \qty{2.433}{\second} in the unperturbed baseline), the extensor rhythm continues through the stimulus without burst suppression, and the post-stimulus period returns to \qty{0.224}{\second} against the \qty{0.221}{\second} baseline. Under the same current delivered at the knee pattern-formation level, the archived project leaves the rhythm-generator phase unchanged ($-\qty{2.3 \pm 2.4}{\percent}$ against baseline), which is the no-effect response that paper reports for its antagonist-only predecessor, the model class this archived project belongs to. The phase-shift magnitudes that paper quantifies belong to its biarticular model, a different network from the archived project, so only the response family is claimed here. This result is a replication in the original engine only, and the program's own port of the same organization has received this current only at an extended stimulus window, not at the paper-exact \SIrange{2}{2.1}{\second} timing.
